\documentclass[12pt]{article}
\usepackage{amsmath, amssymb}
\usepackage[margin=1in]{geometry}
\setlength{\parskip}{6pt}
\title{QUBO Formulation: Facility Location Problem}
\date{}
\begin{document}
\maketitle

\subsection*{Facility Location Problem}

\paragraph{Problem Statement.}
Given a set of $M$ candidate facilities and $N$ customers, the goal is to select a subset of facilities to open and assign each customer to exactly one open facility so as to minimize the total cost. The total cost comprises the fixed opening costs of the selected facilities and the variable assignment costs between customers and their assigned facilities.

\paragraph{Decision Variables.}
Let $y_i \in \{0,1\}$ for $i \in \{0,\dots,M-1\}$ be a binary variable indicating whether candidate facility $i$ is opened ($y_i=1$) or not ($y_i=0$).
Let $x_{ij} \in \{0,1\}$ for $i \in \{0,\dots,M-1\}$ and $j \in \{0,\dots,N-1\}$ be a binary variable indicating whether customer $j$ is assigned to facility $i$ ($x_{ij}=1$) or not ($x_{ij}=0$).
The parameters $f_i \ge 0$ denote the opening cost of facility $i$, and $c_{ij} \ge 0$ denote the assignment cost of customer $j$ to facility $i$.

\paragraph{QUBO Derivation.}
\textbf{Step 1 --- Original Objective.}
The problem is to minimize the total cost:
\begin{align}
\min \quad C(\mathbf{x},\mathbf{y}) = \sum_{i=0}^{M-1} f_i y_i + \sum_{i=0}^{M-1} \sum_{j=0}^{N-1} c_{ij} x_{ij}.
\end{align}

\textbf{Step 2 --- Constraints.}
The feasible region is defined by two sets of constraints:
\begin{align}
\sum_{i=0}^{M-1} x_{ij} = 1 \quad \forall j \in \{0,\dots,N-1\}, \tag{C1}
\end{align}
\begin{align}
x_{ij} \le y_i \quad \forall i \in \{0,\dots,M-1\}, \, j \in \{0,\dots,N-1\}. \tag{C2}
\end{align}

\textbf{Step 3 --- Penalty Term Conversion.}
Introduce a penalty weight $\lambda > 0$. We convert each constraint into a non-negative penalty term that vanishes if and only if the constraint is satisfied.

For (C1), we square the deviation from the required sum:
\begin{align}
P_1 &= \lambda \sum_{j=0}^{N-1} \left( \sum_{i=0}^{M-1} x_{ij} - 1 \right)^2.
\end{align}
Expanding the square and applying the idempotent property $x_{ij}^2 = x_{ij}$ for binary variables:
\begin{align}
\left( \sum_{i=0}^{M-1} x_{ij} - 1 \right)^2 &= \sum_{i=0}^{M-1} x_{ij}^2 + \sum_{i=0}^{M-1} \sum_{\substack{k=0 \\ k \neq i}}^{M-1} x_{ij} x_{kj} - 2 \sum_{i=0}^{M-1} x_{ij} + 1 \nonumber \\
&= \sum_{i=0}^{M-1} x_{ij} + \sum_{i=0}^{M-1} \sum_{\substack{k=0 \\ k \neq i}}^{M-1} x_{ij} x_{kj} - 2 \sum_{i=0}^{M-1} x_{ij} + 1 \nonumber \\
&= 1 - \sum_{i=0}^{M-1} x_{ij} + \sum_{i=0}^{M-1} \sum_{\substack{k=0 \\ k \neq i}}^{M-1} x_{ij} x_{kj}.
\end{align}
Thus, the general contribution of the assignment constraint is:
\begin{align}
P_1 = \lambda \sum_{j=0}^{N-1} \left( 1 - \sum_{i=0}^{M-1} x_{ij} + \sum_{i=0}^{M-1} \sum_{\substack{k=0 \\ k \neq i}}^{M-1} x_{ij} x_{kj} \right).
\end{align}

For (C2), the inequality $x_{ij} \le y_i$ is enforced by penalizing the case $x_{ij}=1$ and $y_i=0$. The corresponding penalty term is:
\begin{align}
P_2 = \lambda \sum_{i=0}^{M-1} \sum_{j=0}^{N-1} (x_{ij} - x_{ij} y_i).
\end{align}
This term is zero when $x_{ij} \le y_i$ and equals $\lambda$ when $x_{ij}=1, y_i=0$.

\textbf{Step 4 --- Combined QUBO Hamiltonian.}
The total energy function $H(\mathbf{x},\mathbf{y})$ is the sum of the objective and penalty terms:
\begin{align}
H = \sum_{i=0}^{M-1} f_i y_i + \sum_{i=0}^{M-1} \sum_{j=0}^{N-1} c_{ij} x_{ij} + P_1 + P_2.
\end{align}
Substituting the expanded forms of $P_1$ and $P_2$:
\begin{align}
H = \sum_{i=0}^{M-1} f_i y_i + \sum_{i,j} c_{ij} x_{ij} + \lambda \sum_{j} \left( 1 - \sum_{i} x_{ij} + \sum_{i \neq k} x_{ij} x_{kj} \right) + \lambda \sum_{i,j} (x_{ij} - x_{ij} y_i).
\end{align}
Collecting linear terms in $x_{ij}$, the $-\lambda x_{ij}$ from $P_1$ and $+\lambda x_{ij}$ from $P_2$ cancel exactly. The Hamiltonian simplifies to:
\begin{align}
H = \sum_{i=0}^{M-1} f_i y_i + \sum_{i=0}^{M-1} \sum_{j=0}^{N-1} c_{ij} x_{ij} + \lambda \sum_{j=0}^{N-1} \sum_{i=0}^{M-1} \sum_{\substack{k=0 \\ k \neq i}}^{M-1} x_{ij} x_{kj} - \lambda \sum_{i=0}^{M-1} \sum_{j=0}^{N-1} x_{ij} y_i + \lambda N.
\end{align}

\textbf{Step 5 --- Q Matrix Entries and Offset.}
Reading off the coefficients from the simplified $H$, the symmetric QUBO matrix entries $Q_{u,v}$ and constant offset $c$ are given in closed form as:
\begin{align}
Q_{y_i, y_i} &= 0, \nonumber \\
Q_{x_{ij}, x_{ij}} &= c_{ij}, \nonumber \\
Q_{x_{ij}, x_{kj}} &= 2\lambda \quad \text{for } i \neq k, \nonumber \\
Q_{x_{ij}, y_i} &= -\lambda, \nonumber \\
Q_{u,v} &= 0 \quad \text{for all other pairs } (u,v), \nonumber \\
c &= \lambda N.
\end{align}

\paragraph{Penalty Weight Justification.}
The penalty weight $\lambda$ must be chosen sufficiently large to ensure that any constraint violation increases the energy more than the maximum possible reduction in the objective function. The maximum possible reduction in the objective from any single variable assignment is bounded by $\max(\max_{i} f_i, \max_{i,j} c_{ij})$. Violating the assignment constraint changes the sum $\sum_i x_{ij}$ from $1$ to $0$ or $2$, incurring a penalty increase of at least $\lambda$. Violating the linking constraint sets $x_{ij}=1$ and $y_i=0$, incurring a penalty of exactly $\lambda$. Therefore, choosing
\begin{align}
\lambda > \max\left( \max_{i} f_i, \max_{i,j} c_{ij} \right)
\end{align}
guarantees that infeasible configurations have strictly higher energy than any feasible configuration.

\paragraph{Correctness Argument.}
Minimizing $H(\mathbf{x},\mathbf{y})$ over $\{0,1\}^{M+MN}$ recovers the optimal solution because the penalty terms are non-negative and vanish exactly on the feasible set. For any feasible assignment, $H(\mathbf{x},\mathbf{y})$ equals the original objective cost $C(\mathbf{x},\mathbf{y})$. For any infeasible assignment, at least one penalty term is strictly positive, adding $\lambda$ (or more) to the energy. Since $\lambda$ exceeds the maximum possible objective value, the energy of any infeasible point strictly exceeds that of the best feasible point. Consequently, the global minimum of $H$ must be feasible and correspond to the configuration that minimizes the original cost.

\paragraph{Illustrative Example.}
Consider a small instance with $M=2$ candidate facilities and $N=2$ customers. The parameters are $f_0, f_1$ for facility opening costs, and $c_{00}, c_{01}, c_{10}, c_{11}$ for assignment costs. The binary variables are mapped to indices $u \in \{0,\dots,5\}$ as follows: $y_0 \to 0$, $y_1 \to 1$, $x_{00} \to 2$, $x_{10} \to 3$, $x_{01} \to 4$, $x_{11} \to 5$. The total number of variables is $n = M + MN = 6$.

The original objective is:
\begin{align}
C = f_0 y_0 + f_1 y_1 + c_{00} x_{00} + c_{10} x_{10} + c_{01} x_{01} + c_{11} x_{11}.
\end{align}
The constraints instantiate to:
\begin{align}
x_{00} + x_{10} = 1, \quad x_{01} + x_{11} = 1, \quad x_{ij} \le y_i \text{ for } i,j \in \{0,1\}.
\end{align}
Applying the general derivation with $\lambda$ as the penalty weight, the expanded penalty terms are:
\begin{align}
P_1 &= \lambda \left[ (x_{00} + x_{10} - 1)^2 + (x_{01} + x_{11} - 1)^2 \right] \nonumber \\
&= 2\lambda - \lambda(x_{00} + x_{10} + x_{01} + x_{11}) + 2\lambda(x_{00}x_{10} + x_{01}x_{11}), \nonumber \\
P_2 &= \lambda \left[ (x_{00} - x_{00}y_0) + (x_{10} - x_{10}y_1) + (x_{01} - x_{01}y_0) + (x_{11} - x_{11}y_1) \right].
\end{align}
Combining objective and penalties, the linear $x_{ij}$ terms cancel, yielding the concrete QUBO Hamiltonian:
\begin{align}
H = &f_0 y_0 + f_1 y_1 + c_{00} x_{00} + c_{10} x_{10} + c_{01} x_{01} + c_{11} x_{11} \nonumber \\
&+ 2\lambda(x_{00}x_{10} + x_{01}x_{11}) - \lambda(x_{00}y_0 + x_{10}y_1 + x_{01}y_0 + x_{11}y_1) + 2\lambda.
\end{align}
The corresponding symmetric Q matrix entries and offset are:
\begin{align}
Q_{0,0} &= 0, & Q_{1,1} &= 0, \nonumber \\
Q_{2,2} &= c_{00}, & Q_{3,3} &= c_{10}, & Q_{4,4} &= c_{01}, & Q_{5,5} &= c_{11}, \nonumber \\
Q_{2,3} &= 2\lambda, & Q_{4,5} &= 2\lambda, \nonumber \\
Q_{0,2} &= -\lambda, & Q_{1,3} &= -\lambda, & Q_{0,4} &= -\lambda, & Q_{1,5} &= -\lambda, \nonumber \\
c &= 2\lambda.
\end{align}
All other entries are zero. Minimizing $H$ over $\{0,1\}^6$ yields the optimal facility selection, customer assignments, and the minimum total cost for this instance.

\end{document}